\begin{figure}[htbp]
\centering
\fbox{\begin{minipage}{0.93\textwidth}
\vspace{2pt}
\raggedright\footnotesize
\textbf{Alfabeto de ação por unidade} --- excerto dos dezoito estados:

\vspace{2pt}
\texttt{GERACAO-PROCEDIMENTO@UNIDADE-A} · \texttt{GERACAO-DOCUMENTO@UNIDADE-A} ·
\texttt{ASSINATURA-DOCUMENTO@UNIDADE-A} · \texttt{CONCLUSAO-AUTOMATICA@UNIDADE-A} ·
\texttt{PROCESSO-REMETIDO@UNIDADE-B} · \texttt{PROCESSO-RECEBIDO@UNIDADE-B} ·
\texttt{GERACAO-DOCUMENTO@UNIDADE-B} · \ldots ·
\texttt{CONCLUSAO-PROCESSO@UNIDADE-A}

\vspace{6pt}
\textbf{Alfabeto de roteamento} --- os mesmos dezoito estados, colapsados em três:

\vspace{2pt}
\texttt{UNIDADE-A} $\rightarrow$ \texttt{UNIDADE-B} $\rightarrow$ \texttt{UNIDADE-A}

\vspace{6pt}
\textbf{Alfabeto de procedimento} --- os mesmos dezoito estados, abstraídos do organograma:

\vspace{2pt}
\texttt{GERACAO-PROCEDIMENTO} $\rightarrow$ \texttt{GERACAO-DOCUMENTO} $\rightarrow$
\texttt{ASSINATURA-DOCUMENTO} $\rightarrow$ \texttt{PROCESSO-REMETIDO} $\rightarrow$
\texttt{PROCESSO-RECEBIDO} $\rightarrow$ \texttt{GERACAO-DOCUMENTO} $\rightarrow$ \ldots
$\rightarrow$ \texttt{CONCLUSAO-PROCESSO}
\vspace{3pt}
\end{minipage}}
\caption{Um mesmo processo administrativo renderizado nos três alfabetos de traço. O alfabeto de
ação por unidade conserva o que se fez e onde; o de roteamento conserva apenas por onde o processo
passou; o de procedimento conserva apenas o que se fez. Os dois últimos colapsam repetições
consecutivas, e é por isso que avaliam populações distintas.}
\label{fig:alfabetos}
\end{figure}